\paragraph{The selected levels and the separating section.}

Equivariance identifies the row as a cusp space. Separating its level components also requires a section that distinguishes those components at the chosen embedding.

The components are those of the complex uniformization of the chosen construction at its fixed complex place. A component $c$ has lower
arithmetic group $\Gamma_c$ and its two prescribed lattices.
The change of level sends $c$ to $hc$ and has a hermitian
transporter $t(h,c)$ carrying both lattices to their counterparts.
The transporter multiplication is well defined modulo the lower group of the destination. For \(h,k\) in the level group it gives
\[
  t(h,kc)t(k,c)t(hk,c)^{-1}\in\Gamma_{hkc}.
\]
 A transporter fixing its
component lies in the lower group only for the identity level element.
The lower character on the arithmetic group is trivial. The
components have the common signature, CM embeddings and weight.
Use the equivariant cusp-form realization of
Section~\ref{subsec:selected-cusp-reading} at this weight. It identifies
the scalar extension of the whole row with the space of bounded
holomorphic sections on these components.
\label{cor:paired-cusp-ranks-1}

Let \(H\) reduce to upper unitriangulars at one chosen level prime \(g\), and let \(H_c\) stabilize the component \(c\). Take the upper group \(G\) generated by \(\Gamma_c\) and the transporters of \(H_c\). The transporter product law above and the trivial kernel on level elements give
\[
  1\longrightarrow\Gamma_c\longrightarrow G\longrightarrow H_c\longrightarrow1.
\]
 It preserves both lattices and is unipotent modulo every prime
of $E$ above $g$. The congruence on the chosen half of the primes
follows from the level change and the kernel congruence for $\Gamma_c$.
On the conjugate half, the reduction of a unitary matrix is the
contragredient with respect to the perfect paired form, so it too
is unipotent.
If \(A\) is unipotent, so is its inverse transpose on the conjugate half, since
\[
  A^{-\mathsf t}-I=-A^{-\mathsf t}(A^{\mathsf t}-I),\qquad (A^{\mathsf t}-I)^r=0 .
\]
 The prime $g$ splits completely in $E$ and exceeds $3r$.
These conditions apply to the character subgroup, the two half-level
subgroups used below, and the trivial subgroup.
\label{cor:paired-cusp-ranks-2}

\begin{lemma}[Separation at the actual two-level translates]
\label{lem:separating-level-section}
Let \(x_b\) denote the coset centers in the chosen cut and \(c\) the selected center. There is a cusp section \(s_c\) of the prescribed determinant weight whose values in the chosen frames satisfy
\[
 s_c(x_b)=\begin{cases}1,&b=c,\\0,&b\ne c.\end{cases}
\]
 The comparison uses the marked subgroups at the two
actual levels and the same coefficient realization as the
character calculation.
\end{lemma}

\begin{proof}
\proofstep{A cusp section distinguishing the level translates.}
Use a balanced majorant of the self-dual component lattice, in a
frame $x_c$ for the Hermitian form and that majorant. The property
needed of the majorant is that an endomorphism of the lattice
strictly contracting at every complex embedding is zero. The
semistable-majorant statement gives this property by applying the
degree inequality to the kernel and image of a contraction. The lattice has degree zero at the semistable majorant, and its
sublattices have degree at most zero. A nonzero strict contraction
would increase the degree of the quotient by its saturated kernel
when mapped isomorphically onto its image lattice. This is
incompatible with the degree bound for that image in the lattice.
Indeed for a nonzero strict contraction \(T\) the saturated kernel gives the incompatible inequalities
\[
  0\le\widehat{\deg}(L/\ker T)<\widehat{\deg}(T L)\le0 .
\]
 The property is invariant under the
right triangular change to an orthonormal frame.

Here is the small-element estimate in this frame. If the operator
norm of $\gamma$ is at most $e^{1/(3r)}$ at every embedding, then
for $1\le a\le r$,
\[
 \operatorname{tr}(\gamma^a)-r\in g\mathcal O_E,
 \qquad |\tau(\operatorname{tr}(\gamma^a)-r)|
       \le r(e^{1/3}+1)<g .
\]
The trace congruence uses nilpotence of the reduction of $\gamma-1$
at each prime over $g$ and squarefreeness of $g\mathcal O_E$. The
field norm now forces all the displayed trace differences to vanish.
Newton's identities give the characteristic polynomial $(X-1)^r$.

In an orthonormal basis the squared Hilbert--Schmidt norm of
$\gamma-1$ is at most $r(e^{2/(3r)}-1)\le e^{2/3}-1<1$, using
$\operatorname{tr}\gamma=r$.
The Hilbert--Schmidt computation uses the trace to cancel both linear terms, giving
\[
  \|\gamma-I\|_{\rm HS}^2=\operatorname{tr}(\gamma^*\gamma)-r\le r(e^{2/(3r)}-1)<1 .
\]
 Its operator norm is strictly below
one, so balance gives $\gamma=1$. The radius potential $R$ on
the positive-plane domain measures the change of majorant from the
orthonormal center. If two points of the ball $R<1/(49r^2)$ are
related by $\gamma$, their two Gram comparisons bound the norm
of $\gamma$ at the center by $e^{1/(3r)}$. This gives a ball whose
nonidentity upper arithmetic translates are disjoint. In the
definite case all upper arithmetic matrices are isometries in the
majorant, so the upper group itself is trivial.
\label{cor:paired-cusp-ranks-3}

For $0<p_0<r$ put $n=f r(r-1)/2$, $d_X=f p_0(r-p_0)$ and
\[
  s_0=2r+50r^2d_X+2,\qquad d=\tfrac12s_0(-1^{[p_0]},+1^{[r-p_0]}).
\]
Here the two signed blocks are repeated at all CM embeddings.
We use the separating section of weight $d$ at the free center.
For clarity, the relevant orbit interpolation estimate concerns an
upper arithmetic group $G$ containing $\Gamma_c$, a free ball
$R<\rho_0=1/(49r^2)$, and the weight potential $\varphi_d$.
The signed two-block calculation used in \hyperref[par:slice-interpolation]{the cut interpolation}
gives $\mathcal L(\varphi_d-\log D)\ge (s_0-2r+2)\mathcal L(R)$.
This margin pays for a radial weight $bR$ and for logarithmic
poles of order $d_X$ at the orbit, since
\[
  s_0-2r+2 > b+d_X/\rho_0.
\]

When \(p_0(r-p_0)\ge4\), use the exponent
\[
  b=4fr^2+1.
\]
 For the smaller signatures
one uses the sum estimate on the positive-plane domain with
$b=4(r-1)$, since the compact flag fibers contribute no factor in
this block-constant weight. In either range the corresponding choice of \(b\) gives a constant \(C\) for the individual model with
\[
  \sup_x\sum_{\gamma\in\Gamma_c}\exp\!\left(-\tfrac b2 R(\gamma x)\right)\le C.
\]
 These bounds use separation in
the orthogonal coordinates and integration of the invariant
volume density over the disjoint balls. All these constants are
for the individual model and weight. Only the stated weight shift
will be used across the family.
\label{cor:paired-cusp-ranks-4}

\proofstep{Interpolation with simultaneous vanishing on the other rows.}
On the flag chart take a smooth cutoff equal to one near the center
and supported inside the free ball. Place smoothed logarithmic
poles at the $G$-orbit of the underlying positive plane. The poles
are truncated to zero outside the disjoint balls. Their Levi loss
is at most $(1+\eta)d_X/\rho_0$ times $\mathcal L(R)$ for a fixed
sufficiently small $\eta>0$.
For the smoothed pole potential \(P\), the curvature cost on these balls is
\[
  {\cal L}(P)\ \ge\ -\frac{(1+\eta)d_X}{\rho_0}\,{\cal L}(R).
\]
 The smoothing parameter tends to zero, with the potential
uniformly bounded on the cutoff's derivative annulus.
Solve $\bar\partial$ of the cutoff with total weight
$\varphi_d-\log D-bR$ plus this logarithmic potential. The same
weighted estimate and frame extension as for the cut interpolation
apply. The solution has uniformly bounded weighted norm as the
smoothings tend to their poles.

In a small ball in the $d_X$ transverse
directions around any orbit point, the weight in the squared-norm integral
is bounded below by a fixed multiple of
$(|z|^2+\varepsilon)^{-d_X}$ times the inverse of the radial
factor. The radial mass diverges logarithmically. For \(d_X>0\) and
a fixed small \(\delta_0>0\), it satisfies
\[
 \int_0^{\delta_0}\frac{r^{2d_X-1}}{(r^2+\varepsilon)^{d_X}}\,dr
    \ \ge\ 2^{-d_X}\int_{\sqrt\varepsilon}^{\delta_0}\frac{dr}{r}
    \ \longrightarrow\ \infty .
\]
Near each orbit centre the cutoff is locally constant, so the
correction there is holomorphic. Its squared sphere mean bounds
its squared value at the centre. The uniform weighted norm bound,
divided by this diverging mass and with the radial weight evaluated
at that centre up to the common comparison constant, makes the
correction at the centre tend to zero. The
bound at each other $G$-orbit point has the additional factor
$e^{-bR/2}$ coming from the radial weight at that point.

Thus after
Poincar\'e summation, the cusp section has readings arbitrarily close
to one at the center and zero at all other coset centers.
For a sequence of such sums \(F_n\), with the representatives indexed by \(b\), the limiting readings have the form
\[
  \lim_{n\to\infty}F_n(x_b)=\begin{cases}1&b=c,\\0&b\ne c.\end{cases}
\]
 Summation
uses the orbit bounds just stated. The bounded section is cuspidal
in this block-constant weight with positive cross-block gap, by the
cusp translation and decay criterion. The map from cusp sections
to the finite tuple of readings has a closed image, being a linear
subspace of a finite-dimensional space. It contains the exact delta
tuple.
Writing \(E\) for the finite tuple evaluation and \(\delta_c\) for the desired delta tuple, closedness yields
\[
  \delta_c\in\overline{\operatorname{im}E}=\operatorname{im}E .
\]
 This gives the separating section on a component, extended by
zero on the other components. At a definite component the upper
group is trivial and the determinant power itself is nonzero there.
\cite{hormander-slice}\label{cor:paired-cusp-ranks-5}

\proofstep{Returning the separating section to the chosen components.}
For each $H$-orbit of components choose a representative $c$ and
its separating section. Sum these sections over the orbits. For the
sum $u_0$, the reading of $g^*u_0$ at $x_c$ vanishes unless $g=1$.
Inside the stabilizer this follows from the separation modulo
$\Gamma_c$, and outside it the component supports are disjoint.
At the transported frame of $kc$ the entry $(h,g)$ of the translate
matrix reads $(hgk)^*u_0$ at $x_c$. Solving the group equation gives the unique nonzero entry of row \(h\):
\[
  hgk=1\quad\Longleftrightarrow\quad g=h^{-1}k^{-1}.
\]
 The determinant therefore has a
nonzero value on every component, including when the group is
noncommutative. This construction is available independently for
all the indicated subgroups with the same weight $d$.
\label{cor:paired-cusp-ranks-6}

\end{proof}

\begin{established}[The flat cut and the determinant-weight gaps]
\label{hyp:three-determinant-weight-ranks}\label{cor:paired-cusp-ranks-7}
At an indefinite lift write \(n=f\binom r2\), \(T_0=T\),
\(L=w_L\), \(F=w_F\), and \(\kappa=F+T_0 L\). On the common
flat cut of \hyperref[par:slice-interpolation]{the cut interpolation},
\(n\ge2\), \(11n^2\le T_0\), and \(s_0=2r+50r^2d_X+2\le T_0\).
Number the flag coordinates \(0,\ldots,r-1\), negative block first,
and put \(k_{\rm vol}(i)=2i-(r-1)\). For
\[
                   U=F-k_{\rm vol}-(n+2)L
\]
the adjacent decreasing gaps inside each block are at least zero,
and \(U_b-U_a\ge1\) for every negative coordinate \(a\) and
positive coordinate \(b\). For \(L\) the corresponding lower
bounds are \(1\) and \(19T_0/10\), respectively.
\end{established}
\begin{proof}
The same light columns and auxiliary section give the flat
intersection for the three weights by
\hyperref[par:slice-interpolation]{the cut interpolation}.
Lemma~\ref{hyp:common-slice-scales} gives \(T_0\ge11n^2\);
\(d_X=fp_0(r-p_0)\le n\), and the polynomial bounds in its proof
give \(2r+50r^2n+2\le T_0\) at large scales.

Set \(d=n+2,\ T'=T_0+d\). For an internal gap \(a\) of the
unsigned row of \(\kappa\), the light gap \(l\) is defined by
\(a+2=T'l+\rho,\ 0\le\rho<T'\).
Lemma~\ref{hyp:height-section-weights} gives \(l\ge1\).
Since \(F=\kappa-T_0L\), after reversing and negating the negative
block the increasing shift in the volume weight contributes two:
\[
 a+2=T'l+\rho,\qquad k_{\rm vol}(i+1)-k_{\rm vol}(i)=2,\qquad
 U_i-U_{i+1}=a-(T_0+d)l+2=\rho\ge0 .
\]

At an unsigned terminal entry write
\(a-d_0=T'l+\rho\), where the two offsets \(d_0\) are \(dn_W+2r\)
and zero. Lemma~\ref{hyp:height-section-weights} gives \(a\ge2T_0^2\)
and \(T'\le41T_0/40\). Consequently
\[
 l>\frac{2T_0^2-T_0}{41T_0/40}-1\ \ge\ \frac{19}{10}T_0
 \quad(T_0\ge 100).
\]
The light rows have nonnegative central shifts of opposite signs,
so every positive coordinate minus every negative coordinate has
at least this gap. The unsigned terminal entries of \(F-dL\) are
\(d_0+\rho\). In the difference of the last positive and the first
negative flag coordinate the central shifts of \(F-dL\) contribute
\(n_P-dn_W\); subtracting \(k_{\rm vol}\) costs \(2(r-1)\).
The resulting difference for \(U\) is at least
\[
 (dn_W+2r)+0+n_P-dn_W-2(r-1)=n_P+2\ge1 .
\]
Internal monotonicity makes this the smallest cross-block difference.
\end{proof}

\begin{proposition}[Dimensions at the three determinant weights]
\label{prop:three-selected-weight-ranks}
Use \(n\ge2\), \(T_0\ge11n^2\), and one common flat cut
of Lemma~\ref{hyp:three-determinant-weight-ranks}, with
\(\delta_c\) points on component \(c\). Put
\(\delta=\sum_c\delta_c\). Then the cusp dimensions satisfy
\[
 \dim V_{\kappa-d}\ge\delta\binom{T_0+n}{n},\qquad
 \dim V_{\kappa}\le\delta\binom{T_0+2n+6}{n},\qquad
 \dim V_{\kappa+2d}\le\delta\binom{T_0+3n+10}{n}.
\]
In particular \(\dim V_{\kappa+2d}\le 2\dim V_{\kappa-d}\). The dimensions used in the comparison
are those of the prescribed cusp spaces.
\end{proposition}

\begin{proof}
\proofstep{The lower rank by interpolation.}
For the lower count, subtract $d$ from $F$. In the interpolation
weight $F-d-k_{\rm vol}-nL$ the cross-block coefficient is at least
$1+3.8T_0-s_0$. This exceeds $4fr^2$, since $fr^2\le4n^2$.
Subtracting one further $L$ leaves the cross-block coefficient
positive, with all within-block coefficients nonnegative. These
are the curvature and additional margin required by the cut
interpolation with zero additional potential. It gives bounded
sections $Q_j$ of weight $F-d$ whose readings at the cut are the
coordinate unit vectors.

For each reading $j$ and $|\alpha|\le T_0$,
form
\[
   Q_j\,w_0^{T_0-|\alpha|}\prod_{a=1}^n S_a^{\alpha_a}.
\]
At least one light factor is cuspidal, including when $|\alpha|=T_0$.
In a linear relation, take the least surviving total degree at the
points and use the local coordinates $S_a/w_0$. Their leading
coefficients read the corresponding $Q_j$ alone. They vanish only
when all the coefficients at that degree vanish, so the family is
independent. For $\delta_c$ cut points the lower dimension at
$\kappa-d$ is $\delta_c\binom{T_0+n}{n}$.
\label{cor:paired-cusp-ranks-8}

\proofstep{The upper rank by injecting into jets.}
For the upper count, the jet map at these same points is injective
through degree $T_0+n+6$. For \(f_0\) of weight \(F+vL\), put
\[
 Q=\frac{\|f_0\|^2\|\theta_2\|^2\|\theta_1\|^2\|\eta_0\|^{2/a}}
        {\big(\sum_{i=1}^n\|S_i\|^2\big)^{v+n+6+1/2}}.
\]
The two theta weights are $(n+4)L-F$ and $2L-d_{\rm sign}$,
and the auxiliary weight is $a d_{\rm sign}$ for the determinant
sign weight $d_{\rm sign}=(-1^{[p_0]},+1^{[r-p_0]})$.
The effective numerator weight is
\[
 (F+vL)+((n+4)L-F)+(2L-d_{\rm sign})+d_{\rm sign}
               =(v+n+6)L .
\]
At a cut point write \(\varrho\) for the size of the coordinate
tuple \(S_i/w_0\).

If the jets of \(f_0\) through degree \(v+n+6\) vanish, its
squared norm vanishes to order at least \(2(v+n+7)\). The cut rows are
local coordinates, so their sum of squared norms is comparable to
\(\varrho^2\). In this frame,
\[
 Q=O\!\left(\varrho^{\,2(v+n+7)-2(v+n+6+1/2)}\right)
   =O(\varrho)\qquad(\varrho\to0).
\]
At a noncompact end the lower bound for the light sum and the
numerator decay give
\[
  Q\ll 
       \exp\bigl((-4/g_0+2c(v+n+6+1/2)) P\bigr).
\]
For \(v=T_0\) the coefficient of \(P\) is negative by
\(T_0cg_0\le1\) and \(n+6+1/2<T_0\). A positive maximum could only occur
where the factors in the numerator are nonzero and the light sum
is positive.

At such a point the logarithm of the quotient is a
sum of pluriharmonic terms, a negative multiple of
$\log\sum_a|S_a|^2$, and $\varphi_L/2$. The numerator contributes zero Levi form at this point.
The map \(\lambda:z\mapsto[S_1(z):\cdots:S_n(z)]\) has target
dimension \(n-1\) and therefore has a nonzero tangent vector
\(\xi\) in the kernel of its differential. Along \(\xi\) the
pullback of the Fubini--Study form vanishes, so
\[
   \mathcal L(\log\text{quotient})(\xi,\bar\xi)
        =\tfrac12\mathcal L(\varphi_L)(\xi,\bar\xi)>0 .
\]
The last inequality uses the positive curvature of the selected
light weight on the flag chart. A smooth function with a local
maximum has nonpositive Levi form there. Thus all-zero jets force
\(f_0=0\). The dimension at $\kappa$ is
at most $\delta_c\binom{T_0+2n+6}{n}$.

\proofstep{The second upper weight shift.}
The cut contains a point on each component under consideration.
The bounded base theta section of weight $F$ has a nonzero value
at some frame. At that frame the nonvanishing light family gives
\[
  |w_0|^2+\sum_{i=1}^n|S_i|^2>0.
\]
 Multiply the base section by its
$T_0$th power to obtain a nonzero bounded section of weight
$F+T_0L$. The jet injection just proved applies to that bounded
section and would force it to vanish if the cut were empty.

For the upper count at $\kappa+2d$, the weight $(n+4)L-2d$
has positive within-block and cross-block differences. In the
curvature test subtract both $nL$ and the density weight. After both subtractions the remaining cross-block margin satisfies
\[
  7.6T_0-2s_0-2(r-1)>4fr^2.
\]
 Subtracting a further $L$
leaves nonnegative differences. These bounds give bounded
section interpolation on the same cut. Multiplication by one such nonzero
bounded section is injective by holomorphic continuation, and
puts the result at $F+(T_0+n+4)L$.

The maximum argument applies
there with jet degree $T_0+2n+10$. The extra $n+4$ in the
end estimate uses $2n+10.5<T_0$ and $T_0cg_0\le1$, so
$2c(v+n+6.5)<4/g_0$ for $v=T_0+n+4$. Hence the upper count at $\kappa+2d$ is
$\delta_c\binom{T_0+3n+10}{n}$. The component upper bounds
sum by restriction, and the lower bound sums by extension by zero.

For the upper shift we have
\[
 \frac{\binom{T_0+3n+10}{n}}{\binom{T_0+n}{n}}
 \le \exp\left(\frac{n(2n+10)}{T_0+n}\right)<2
 \quad(n\ge2,\ T_0\ge 11n^2).
\]
The loss in the lower shift relative to the unshifted count is
used below. The numbers of cut points and the interpolation
constants disappear from the ratios.
\label{cor:paired-cusp-ranks-9}
\end{proof}
